\honest{Decoherence is Falsifiable and Testable}{Eliezer Yudkowsky}{2008}

People use ``falsifiable'' and ``testable'' as if they meant the same thing. I use them for two different things. \nb{Popper used them as synonyms: ``Testability is falsifiability.'' The split is the post's own, and it says so.}

I start with Bayes's theorem, which I own on at least two items of clothing. A hypothesis is ``falsifiable,'' in my sense, if it concentrates its probability on some outcomes, so that others would drive its probability nearly to zero. ``Thus is the rule of science derived in probability theory.'' \nb{A tiny likelihood gives a tiny posterior only if some rival explains the outcome better; the post grants this a few paragraphs later. Popper himself held, in Barnes's summary, that ``it was impossible to show that a theory was even probable based on the evidence.''}

Is decoherence falsifiable? \nb{As in the previous post, ``decoherence'' here means the many-worlds view. Physicists also use the word for a process, the suppression of interference by the environment, which experiments observe and which does not by itself decide between interpretations.} Yes: entangled particles that should have opposite spins could turn up with the same spin, or apples could fall upward. An imagined objector says that this only tests quantum mechanics, not decoherence as against collapse, and I agree. Popper's complaint was of this single-hypothesis kind: psychoanalysis could explain anything a patient did. \nb{That is accurate. Popper called the psychoanalytic theories ``simply non-testable, irrefutable.''}

A being who knew only the wave equation ``plus the Born probability rule'' would have sharp expectations. So would someone who knew the wave equation plus a collapse postulate. Both theories pass the falsifiability test.

Then ``testable.'' A prediction is ``new'' only relative to another hypothesis. A test is an outcome that has different probabilities under decoherence and under collapse, and Bayes's theorem treats the two hypotheses alike. If decoherence is untestable relative to collapse, collapse is untestable relative to decoherence. Had Everett and Wheeler come before Bohr and Heisenberg, would we demand new predictions from collapse? A reasoner whose beliefs depend on the order of the evidence suffers from ``hysteresis,'' which I call a Bad Thing. \nb{The demand for new predictions has a name, predictivism, and Popper is its best-known proponent. It has Bayesian defenders too, such as Patrick Maher, for whom a successful prediction is extra evidence that the method that produced the theory is reliable. The post does not mention them.}

What about crackpots who add untestable parts to working theories, such as angels who push charged particles? Reject them, I agree, but not because their ideas are new. Suppose a Church-funded science had read Maxwell's laws as the laws that direct angels. A newcomer who proposed to keep the equations and drop the angels would be told that the proposal makes no new predictions. ``It seems to me'' that the old angels should go for the same reason as the new ones: they are useless. A Bayesian charges each theory for the length of its description, and then the order in which you hear the theories does not matter.

The objector replies that the many worlds are exactly such a useless addition. No: what decoherence asks you to believe is the wave equation, and ``just look at a hydrogen atom'' to test it. \nb{The collapse view keeps the same equation and adds a collapse, so these tests confirm what the two views share.} The worlds are ``just a deductive consequence of the wavefunction's evolution,'' as the spaceship beyond the cosmological horizon follows from conservation laws. There are critiques, ``most notably'' the question of where the Born probabilities come from. \nb{Six days earlier, in ``The Born Probabilities,'' Yudkowsky had called this ``One serious mystery of decoherence.''} To deny the worlds is ``necessarily'' to deny that the quantum laws hold everywhere. \nb{Bohmian mechanics keeps the wave equation everywhere, with no collapse, and has one world; see ``Many Worlds, One Best Guess.''}

This controversy ``may be unique in scientific history'': serious scientists are saying that a successful theory contains angels. \nb{In 1905 Einstein called the ether ``superfluous,'' and Lorentz's ether theory, by then empirically equivalent to special relativity, lost out for that reason. The precedent supports the post's rule.} The angel is the collapse postulate, which can be thrown away, leaving the theory ``strictly simpler.'' \nb{In the textbook theory the collapse postulate supplies the probabilities, and the post's own decoherence adds them back as ``the Born probability rule.'' Whether many-worlds is then simpler is disputed among philosophers of physics. The book returns to the question in ``Many Worlds, One Best Guess.''}

My conclusion: decoherence ``is not ruled out by Occam's Razor, nor is it unfalsifiable, nor is it untestable.'' \nb{For ``untestable,'' the post has argued only that untestability would count against both theories; it names no outcome that separates them. Such outcomes exist in principle, which supports the title: Vaidman writes that the claim that many-worlds cannot be told apart from an ideal collapse theory ``is not the case,'' and experiments have ruled out some physical collapse models. ``Collapse Postulates,'' placed two posts earlier in the book, makes this point.} The two theories should be compared on their fit with experience, and the arguments that keep decoherence out of the game are, ``I would argue,'' an outright error of probability theory.
